\section*{Epilogue}
\addcontentsline{toc}{section}{Epilogue}
\label{sec:a-epilogue}

一份回顾式讲义如果只讲已经建好的部分，读完的感觉会比实际状况好太多。所以最后一节专门列缺口。
下面每一条都当场核过——用 \code{ls}、\code{grep} 或者读文件本身——核的办法写在括号里，
方便你几个月后再验一遍。

\subsection*{缺口清单}

\begin{table}[ht]
\centering
\begin{tabularx}{\textwidth}{@{}>{\raggedright\arraybackslash}X>{\raggedright\arraybackslash}X@{}}
\toprule
缺什么 & 怎么核实的 \\
\midrule
Week 3 一行代码都没有 &
\codet{inferenceLM/} 下 \codeb{grep -rniE "batch|schedul|admission|drain|shutdown"} 只命中三处，
全是 \codeb{transformer_lm_adapter.py} 注释里描述 tensor shape 的 \codet{(batch, seq)}。
\codet{git log} 停在 \codet{Week3 start}（2026-04-21） \\
\addlinespace
\codet{inferenceLM/output/} 是空的 &
目录里只有一个 0 字节的 \codet{\_\_init\_\_.py}。detokenize 与流式回传在
\codet{DESIGN.md} §3 的架构图里画着，从没落地。今天调用方只能轮询
\codeb{request_store[request_id]} \\
\addlinespace
\codet{benchmarks/} 是空的 &
目录里只有一个 \codet{.gitkeep}。这个仓库里没有任何性能测量代码；
唯一一个真实穿过 \codet{LMEngine} 的 serving benchmark 在姊妹项目里，不在这里 \\
\addlinespace
没有 HTTP 层 &
\codeb{grep -rniE "socket|fastapi|uvicorn|http|flask|grpc"} 在 \codet{inferenceLM/} 下零命中。
系统的入口 API 就是 \codeb{RequestReceiver.submit_request()} 这个协程；
\codet{DESIGN.md} §3 画的 "TCP/stdin" 从未存在 \\
\addlinespace
没有 CI &
没有 \codet{.github/} 目录。\codet{ruff} 和 \codet{mypy} 声明在
\codet{pyproject.toml} 的 dev 依赖里（第 18--22 行），但全文件没有
\codet{[tool.ruff]} 或 \codet{[tool.mypy]} 段，也没有任何东西调用它们 \\
\addlinespace
\codet{LICENSE} 是 0 字节 &
\codet{ls -la} 显示 size 为 0。文件在，内容没有 \\
\addlinespace
\codet{README.md} 还是脚手架 &
第 2 行标题写着 \codeb{[Project Name --- TO BE DECIDED]}，第 7 行状态写着
"Setup phase --- project topic and tech stack under discussion"。
全文 30 行，多数是目录树注释 \\
\addlinespace
两个状态值是死的 &
\codet{RequestStatus.PREFILLING} 和 \codet{DECODING} 定义在
\codeat{inferenceLM/data/request\_status.py}{5}--\codeat{inferenceLM/data/request\_status.py}{6}，
但 \codet{grep} 全仓库，除了枚举定义本身和第 10 行那句注释之外，没有任何一处引用 \\
\addlinespace
\codet{pytest-asyncio} 装错了位置 &
它在 \codeat{pyproject.toml}{11}，属于运行时 \codet{dependencies}，
而不是 \codet{[dependency-groups] dev}。装这个包的人不一定要跑测试 \\
\bottomrule
\end{tabularx}
\caption{核实过的缺口清单。前三条是能力缺失，后面几条是收尾没做完；
两类的修复成本差着一个数量级，但都在同一份 \codet{ROADMAP.md} 里排着队。}
\label{tab:epilogue-gaps}
\end{table}

\subsection*{先做哪三个}

排序不用重排——仓库根的 \code{ROADMAP.md} Part I 已经排好了，依据是\cnemph{依赖关系而不是价值大小}。
这里只把最前面三条抄过来，各配一句"第一步动什么"。

\textbf{第一，engine 生命周期状态机（\code{INIT} $\to$ \code{RUNNING} $\to$ \code{DRAINING}
$\to$ \code{STOPPED}）。}
它是唯一真正的阻塞项，也是 \code{BACKLOG.md} 里仅剩的有效 P0。
今天 \code{kill()} 只把 \code{self.open} 置 False，而循环正卡在
\code{await pending_queue.get()} 上——空队列上的一个 killed engine 会永远挂着。
第 4 章那条"测不到关闭路径"就是这件事的症状。
\cnemph{第一步：}给 \code{InferenceEngine} 加显式状态字段和 \codeb{async shutdown(drain=True)}，
然后把测试里的 \code{task.cancel()} 一个一个换成它。换不掉的那个，就是设计还没想清楚的地方。

\textbf{第二，把 request lifecycle 和 iteration lifecycle 拆开。}
这是 \secref{sec:a-batching} 全章的根因。
\cnemph{第一步：}把 \codeb{_keep_get_request_and_inference()} 从"取一条跑到底"
改成"每轮迭代推进所有在飞请求各一步"，为此需要一个 in-flight 集合来存
\codeb{(request_id, latest_token, past_key_values)}。
动手之前先解决一个文档冲突：\code{future-tasks.md} 分了 Scheduler 和 Batcher 两层，
而 \code{week3-discussion.md} 明说不做 pluggable scheduler——先定哪份是现行意图，再写代码。

\textbf{第三，stage-homogeneous batching。}
已经 LOCKED，理由是后端约束不是偏好（\secref{sec:cb-stage}）。
\cnemph{第一步：}每轮迭代改成两次 \code{model()} 调用，先 prefill microbatch 再 decode microbatch，
并且从第一行代码起就记住这两次是\cnemph{顺序执行}的。
再往后是 admission policy——它在 \code{ROADMAP.md} 里排第四，
并且是整条路线上唯一一个还没有答案的设计问题；建议等前两步的循环真能跑起来、
拿到真实的 TTFT/ITL 曲线之后再定，而不是反过来。

\subsection*{这个仓库到底是什么}

数一下行数：\code{inferenceLM/} 下 13 个 Python 文件共 \cnemph{514} 行（其中 4 个 \code{__init__.py}
是 0 行），\code{tests/} 共 \cnemph{880} 行，\code{docs/} 下自己写的 markdown 共 \cnemph{1502} 行
（不含四份模板）。文档大约是源码的\cnemph{三倍}，而测试比源码还多。

一个五百行的 Python 包，跑得动一条请求的完整 prefill$\to$decode，
支持一种采样策略（greedy），没有 HTTP 层，没有 batching，没有性能数据。
按功能清单算，这份东西很薄。

但清单不是它的产物。真正留下来的是那些\cnemph{被否决的方案}，
以及每一次否决时写下的理由：
为什么不用 \code{model.generate()}（ADR-001，Status 至今还是 Proposed）；
为什么 \code{run()} 必须返回 \code{Task} 而不能 \code{await}；
为什么不提前暴露 \code{do_sample}；
为什么"GPU 利用率低"是症状不是根因；
为什么 \code{DESIGN.md} 里那套 75.5\,MB/slot 的预算表算得没错却已经失效；
为什么 2:1 这个比例不是被换掉了而是整类规则失效了；
以及最后那条——\code{week3-discussion.md} 的 §3 目标清单里，
那条已经在 §2.4 和 §4 被宣告死亡的 2:1 规则，至今还原样留着。

功能可以在一个下午被别的库替代掉。上面这些判断不能。
如果这份讲义只值得记住一件事，那就是：\cnemph{一个设计文档最有价值的部分，
是它记下的那些"当时为什么没这么做"，以及它没能自我更新的那几行。}
